\documentclass[aspectratio=169,10pt,spanish]{beamer}

\input{../../../_preambulo_beamer.tex}
\input{../../../_comandos_beamer.tex}

\selectlanguage{spanish}

\title{MA1001B - Modelación Estadística para la Toma de Decisiones}
\subtitle{Lección 23: Distribución muestral de la media}
\author{}
\institute{Tecnol\'ogico de Monterrey}
\date{}

\begin{document}

\begin{frame}[plain]
  \vspace{1.2cm}
  {\Large\bfseries MA1001B - Modelación Estadística para la Toma de Decisiones\par}
  \vspace{0.7cm}
  {\large Lección 23: Distribución muestral de la media\par}
  \vspace{0.7cm}
  {\color{TecRojo}\rule{\textwidth}{1.2pt}}
  \vspace{0.8cm}

  Facultad MA1001B\\[0.3cm]
  Periodo\\[0.3cm]
  Tecnol\'ogico de Monterrey
\end{frame}

\begin{frame}{La pregunta de la media muestral}
  \begin{alertblock}{Pregunta guía}
    Si los tiempos de ciclo son $N(80,12)$, ¿cuánto se mueve $\bar x$ cuando
    $n=9$ frente a $n=36$?
  \end{alertblock}

  \vspace{0.6em}
  Al final de esta lección, usted puede:
  \begin{itemize}
    \item escribir $E[\bar x]=\mu$ y $\mathrm{EE}(\bar x)=\sigma/\sqrt{n}$;
    \item comparar distribuciones muestrales para dos tamaños de muestra;
    \item verificar la fórmula del EE con una simulación sembrada;
    \item explicar por qué las copias agrupadas inflan el EE observado.
  \end{itemize}

  \vfill
  \scriptsize Todos los tiempos de ciclo usados hoy son sintéticos. No se usan datos reales de empresa.
\end{frame}

\begin{frame}{El error estándar se reduce con $n$}
  \small
  Población $X\sim N(80,12^2)$. Para una muestra iid,
  \[
    E[\bar x]=80,\qquad
    \mathrm{EE}(\bar x)=\frac{12}{\sqrt{n}}.
  \]
  Entonces
  \[
    \mathrm{EE}(9)=\frac{12}{3}=4,\qquad
    \mathrm{EE}(36)=\frac{12}{6}=2.
  \]
  Cuadruplicar $n$ reduce el EE a la mitad.

  \vspace{0.3em}
  \begin{exampleblock}{Población sintética de tiempos de ciclo}
    $\sigma=12$ describe una observación. $\sigma/\sqrt{n}$ describe $\bar x$.
  \end{exampleblock}
\end{frame}

\begin{frame}[fragile]{Paso 1: Fórmula del error estándar}
  \begin{columns}[T]
    \begin{column}{0.42\textwidth}
      \small
      \begin{lstlisting}
se_9 = 12 / sqrt(9)
se_36 = 12 / sqrt(36)
      \end{lstlisting}

      \begin{block}{Salida verificada}
        EE$(n=9)=4.000000$\\
        EE$(n=36)=2.000000$\\
        Razón $=2.000000$
      \end{block}
    \end{column}
    \begin{column}{0.58\textwidth}
      \centering
      \includegraphics[width=\textwidth]{../figuras/sampling_mean_01_standard_error.png}
      \vspace{0.2em}
      \scriptsize Densidades poblacional y muestrales del Paso 1
    \end{column}
  \end{columns}
\end{frame}

\begin{frame}{La simulación recupera la fórmula del EE}
  \small
  Extraer 5{,}000 muestras iid con semilla 42.
  \[
    n=9:\ \text{media de las medias}=80.036617,\ \widehat{\mathrm{EE}}=3.997160.
  \]
  \[
    n=36:\ \text{media de las medias}=79.994341,\ \widehat{\mathrm{EE}}=2.028859.
  \]
  Ambos EE empíricos quedan junto a las fórmulas 4 y 2.
\end{frame}

\begin{frame}[fragile]{Paso 2: $n=9$ versus $n=36$}
  \begin{columns}[T]
    \begin{column}{0.40\textwidth}
      \small
      \begin{lstlisting}
draws = rng.normal(
  80, 12, size=(5000, n)
)
xbars = draws.mean(axis=1)
      \end{lstlisting}

      \begin{block}{Salida verificada}
        $n=9$ EE $=3.997160$\\
        $n=36$ EE $=2.028859$
      \end{block}
    \end{column}
    \begin{column}{0.60\textwidth}
      \centering
      \includegraphics[width=\textwidth]{../figuras/sampling_mean_02_n9_vs_n36.png}
      \vspace{0.2em}
      \scriptsize Distribuciones muestrales simuladas del Paso 2
    \end{column}
  \end{columns}
\end{frame}

\begin{frame}{Paso 3: El EE supone extracciones iid}
  \begin{columns}[T]
    \begin{column}{0.42\textwidth}
      \small
      Cada ``muestra'' de 36 son 6 valores de conglomerado copiados 6 veces.

      \vspace{0.4em}
      El EE de fórmula para $n=36$ es $2$.\\
      EE empírico agrupado $=4.992583$.\\
      EE de fórmula para 6 conglomerados es $4.898979$.

      \vspace{0.4em}
      \textbf{Límite:} el EE supone extracciones iid de la población definida.
      La Lección 24 invoca el TLC para un reloj asimétrico.
    \end{column}
    \begin{column}{0.58\textwidth}
      \centering
      \includegraphics[width=\textwidth]{../figuras/sampling_mean_03_iid_limit.png}
      \vspace{0.2em}
      \scriptsize Medias iid versus agrupadas del Paso 3
    \end{column}
  \end{columns}
\end{frame}

\begin{frame}{Verificación de conceptos}
  \begin{enumerate}
    \item Calcule $\mathrm{EE}(\bar x)$ para $n=9$ cuando $\sigma=12$.
    \item ¿Por qué cuadruplicar $n$ de 9 a 36 reduce el EE a la mitad?
    \item El EE empírico de $n=36$ es $2.028859$. ¿Sustituye eso a $2$?
    \item ¿Por qué una muestra de 36 copias agrupadas se acerca a $n=6$?
  \end{enumerate}

  \vspace{0.6em}
  \begin{block}{Conexión con la Lección 24}
    La sesión puede continuar directamente con el teorema del límite central
    para una población sesgada a la derecha.
  \end{block}

  \vfill
  \scriptsize Fuente: OpenStax, \textit{Introductory Business Statistics 2e},
  Capítulo 7, Sección 7.1. CC BY-NC-SA.
\end{frame}

\appendix



\end{document}
